% ===================================================================
%  TABLEAU N° 13 — L'Hôtel. --- Le Restaurant. --- Le Café.
%  Feuillets 75 a 77 du fac-simile = folios 73 a 75.
%  Correspondance : folio imprime = numero de feuillet - 2.
%
%  Les cles « %%K » sont les MEMES que dans texto/io/22-tabelo-13.tex :
%  c'est par elles que la page de lecture met les deux textes en
%  regard. Ce tableau porte DEUX notes de bas de page numerotees
%  « (1) », l'une folio 73 (detail de la chambre, renvoi au tableau
%  6 -- cf. la note io correspondante, en fac-simile io folio 85), et
%  l'autre folio 74 (le menu suit la liste des mets du vocabulaire --
%  cf. la note io, en fac-simile io folio 87). La pagination des deux
%  notes ne coincide donc PAS entre les deux editions : ecart reel de
%  mise en page, non une erreur de releve.
%
%  ATTENTION — LE FAC-SIMILE FRANCAIS EST UNE PRISE DE VUE, contraste
%  inegal d'une page a l'autre ; plusieurs renvois chiffres a demi
%  effaces (conducteur, malle, voyageur, hotelier, telegraphiste,
%  depeche, touriste, appareil photographique, guide, Ocean, boite aux
%  lettres) ont ete recoupes avec la legende numerique de
%  texto/io/22-tabelo-13.tex, qui porte les memes numeros d'objets.
% ===================================================================

% --- Feuillet 75 (folio 73, ouverture de tableau) -------------------
\begin{VUpage}[75]{73}
\VUnotes{143.2mm}{%
\VUgras{(1) Voir,} pour les détails de la chambre, le tableau n\textsuperscript{o} 6.%
}

%%K t13-apar-1 apar
\VUsaut{3.0mm}
\VUtitre{15.0pt}{60}{TABLEAU N\textsuperscript{o} 13}
\VUsaut{1.4mm}
\VUtitre{11.4pt}{0}{L'Hôtel. --- Le Restaurant. --- Le Café.}
\VUsaut{2.2mm}
\VUfilet{20mm}
\VUsaut{6.4mm}

%%K t13-01-1 p
\VUblancAlinea
1. --- Arrivé hier soir à Royan, je suis descendu\nl
à \textit{l'Hôtel de l'Océan}, qui m'avait été recommandé par\nl
un ami.

%%K t13-01-2 p
\VUblancAlinea
On m'a donné une jolie \VUgras{chambre} \textsuperscript{(2)}, au \VUgras{deuxième}\nl
\VUgras{étage} \textsuperscript{(1)}, où j'ai très bien dormi (1).

%%K t13-02-1 p
\VUblancAlinea
2. --- Ce matin, en sortant de ma chambre, où la\nl
\VUgras{bonne} \textsuperscript{(3)} m'avait apporté mon premier déjeuner, je\nl
suis allé au \VUgras{fumoir} \textsuperscript{(5)} qui, en même temps, sert de\nl
\VUgras{cabinet de lecture} \textsuperscript{(4)}, pour parcourir les \VUgras{jour}\cc
\VUgras{naux} \textsuperscript{(6)}, en fumant une \VUgras{cigarette} \textsuperscript{(8)}. Ma lecture\nl
finie, je suis descendu par l'\VUgras{ascenseur} \textsuperscript{(9)} et je suis\nl
allé me promener en ville.

%%K t13-03-1 p
\VUblancAlinea
3. --- Comme j'avais oublié de demander au \VUgras{bu}\cc
\VUgras{reau} \textsuperscript{(12)} de l'hôtel à quelle heure on servait le repas\nl
de la \VUgras{table d'hôte} \textsuperscript{(11)}, je suis rentré de bonne heure\nl
et j'en ai profité pour aller prendre un bain dans la\nl
\VUgras{salle de bain} \textsuperscript{(10)} de l'hôtel. Je suis ensuite revenu à\nl
la \VUgras{caisse} \textsuperscript{(15)} pour téléphoner à un de mes amis. Mais\nl
le \VUgras{téléphone} \textsuperscript{(13)} était occupé; voyant mon embarras,\nl
la \VUgras{caissière}\textsuperscript{(14)}, qui relevait une \VUgras{note} \textsuperscript{(17)} sur le \VUgras{livre}\nl
\VUgras{des voyageurs} \textsuperscript{(16)}, pria le \VUgras{concierge}\textsuperscript{(18)} d'envoyer\nl
le \VUgras{chasseur} \textsuperscript{(19)} faire ma commission avec sa \VUgras{bicy}\cc
\VUgras{clette} \textsuperscript{(20)}.

%%K t13-04-1 p
\VUblancAlinea
4. --- Je la remerciai et je sortis pour donner un\nl
pourboire à l'\VUgras{homme de peine} \textsuperscript{(24)} qui apportait de\nl
la gare, dans sa \VUgras{voiture à bras} \textsuperscript{(25)}, mon \VUgras{carton à}\nl
\VUgras{chapeau} \textsuperscript{(26)}, ma \VUgras{valise} \textsuperscript{(27)} et mon \VUgras{sac de voyage} \textsuperscript{(28)}.\nl
Le \VUgras{facteur} \textsuperscript{(21)}, qui arrivait à ce moment, me remit\nl
une \VUgras{lettre} \textsuperscript{(22)}.

%%K t13-05-1 p
\VUblancAlinea
5. --- Je restai encore quelque temps sur le seuil de\nl
la porte, près de la \VUgras{boîte aux lettres} \textsuperscript{(23)}, à regarder
\parplein
\end{VUpage}

% --- Feuillet 76 (folio 74) ------------------------------------------
\begin{VUpage}[76]{74}
\VUnotes{143.2mm}{%
(1) \VUgras{\textsc{Nota.}} --- A l'aide de la liste des mets contenue dans le\nl
vocabulaire, le professeur pourra varier facilement le menu ci-dessous.%
}

%%K t13-05-1 p suite
\VUcontinue
le mouvement de la rue. Le \VUgras{conducteur} \textsuperscript{(30)} de\nl
l'\VUgras{omnibus} \textsuperscript{(29)} chargeait sur sa voiture la \VUgras{malle}\nl
d'un \VUgras{voyageur} \textsuperscript{(32)} auquel l'\VUgras{hôtelier} \textsuperscript{(33)} faisait ses\nl
adieux. Un petit \VUgras{télégraphiste} \textsuperscript{(35)} apportait, sans\nl
se presser, une \VUgras{dépêche} \textsuperscript{(34)} pour un client de l'hôtel;\nl
un \VUgras{touriste} \textsuperscript{(36)}, son \VUgras{appareil photographique} \textsuperscript{(38)}\nl
en bandoulière, consultait son \VUgras{guide} \textsuperscript{(37)}.

%%K t13-06-1 p
\VUblancAlinea
6. --- Devant la grille, un placide \VUgras{commission}\cc
\VUgras{naire} \textsuperscript{(39)} sommeillait entre ses \VUgras{crochets} \textsuperscript{(40)} et sa\nl
\VUgras{boîte à cirer} \textsuperscript{(41)}, tandis qu'une petite \VUgras{blanchis}\cc
\VUgras{seuse} \textsuperscript{(42)}, qui portait un grand \VUgras{panier} \textsuperscript{(43)} de linge,\nl
riait de l'air solennel du \VUgras{maître d'hôtel} \textsuperscript{(44)}, qui se\nl
tenait auprès de la porte.

%%K t13-07-1 p
\VUblancAlinea
7. --- La vue du \VUgras{cuisinier} \textsuperscript{(45)}, qui était sorti de\nl
sa \VUgras{cuisine} \textsuperscript{(47)}, située dans le \VUgras{sous-sol} \textsuperscript{(48)}, pour\nl
envoyer un \VUgras{marmiton} \textsuperscript{(46)} faire une commission, me\nl
rappela que l'heure du déjeuner était proche et je me\nl
rendis au restaurant en traversant la cour, dans\nl
laquelle un \VUgras{chauffeur} \textsuperscript{(50)} garait son \VUgras{automobile} \textsuperscript{(49)},\nl
tandis qu'un \VUgras{mécanicien} \textsuperscript{(52)} réparait, sur les indica\cc
tions du \VUgras{cycliste} \textsuperscript{(53)}, un \VUgras{tricycle à pétrole} \textsuperscript{(51)},\nl
auquel il venait d'arriver un accident.

%%K t13-08-1 p
\VUblancAlinea
8. --- Je demandai à un petit \VUgras{groom} \textsuperscript{(54)} qui portait\nl
des \VUgras{bocks} \textsuperscript{(55)}, si la table d'hôte était sous la véranda,\nl
et sur sa réponse affirmative, je pris place à l'une des\nl
tables et je me fis servir un excellent repas.

%%K t13-09-1 p
\VUblancAlinea
9. --- Le garçon m'apporta le menu (1) du déjeuner\nl
et la carte des vins. Je fis mon choix et je lui com\cc
mandai en \VUgras{hors-d'œuvre} des \VUgras{crevettes} et du\nl
\VUgras{beurre}, et, pour entrée, des \VUgras{tripes à la mode de}\nl
\VUgras{Caen}. Je ne pris pas de \VUgras{poisson}, mais je demandai\nl
une portion de \VUgras{gigot} d'agneau et, comme \VUgras{légume},\nl
des \VUgras{épinards} à la crème. Puis, aucun des \VUgras{entremets}
\parplein
\end{VUpage}

% --- Feuillet 77 (folio 75) ------------------------------------------
\begin{VUpage}[77]{75}
%%K t13-09-1 p suite
\VUcontinue
ne me plaisant, je me fis apporter des \VUgras{desserts}\nl
variés, \VUgras{fromage}, \VUgras{fruits} et \VUgras{biscuits}.

%%K t13-09-2 p
\VUblancAlinea
J'avais arrosé ce menu d'un excellent \VUgras{vin} de Saint\cc
Émilion, et je quittai la table très satisfait de mon\nl
repas, après avoir gratifié le \VUgras{garçon} \textsuperscript{(57)} d'un bon\nl
\VUgras{pourboire} \textsuperscript{(59)}.

%%K t13-10-1 p
\VUblancAlinea
10. --- En me voyant prêt à partir, le \VUgras{gérant} \textsuperscript{(60)}\nl
me proposa de passer sur le balcon pour admirer le\nl
superbe panorama de l'\VUgras{Océan} \textsuperscript{(62)}. On apercevait au\nl
loin de petites \VUgras{barques} \textsuperscript{(63)} de pêche, des \VUgras{voiliers} \textsuperscript{(64)}\nl
et un énorme \VUgras{paquebot} \textsuperscript{(65)}, dont la silhouette noire\nl
se détachait sur les \VUgras{nuages} \textsuperscript{(67)} blancs de l'\VUgras{horizon} \textsuperscript{(66)}.\nl
Une brise légère faisait flotter le \VUgras{drapeau} \textsuperscript{(70)} planté\nl
au-dessus de l'\VUgras{enseigne} \textsuperscript{(69)} du \VUgras{hall} \textsuperscript{(68)}.

%%K t13-11-1 p
\VUblancAlinea
11. --- Le spectacle était si intéressant que je me\nl
promis de monter le lendemain sur la terrasse du\nl
café en emportant une \VUgras{lorgnette} \textsuperscript{(71)} ou une \VUgras{longue}\cc
\VUgras{vue} \textsuperscript{(72)}.

%%K t13-12-1 p
\VUblancAlinea
12. --- En quittant le balcon, je me rendis au café\nl
de l'hôtel pour prendre une \VUgras{consommation} \textsuperscript{(73)}, en\nl
faisant une partie de \VUgras{billard} \textsuperscript{(75)}. Je traversai sans\nl
m'arrêter la salle du café où de nombreux consom\cc
mateurs jouaient aux \VUgras{échecs} \textsuperscript{(78)}, aux \VUgras{dames} \textsuperscript{(79)},\nl
aux \VUgras{dominos} \textsuperscript{(80)}, au \VUgras{trictrac} \textsuperscript{(81)} ou aux \VUgras{cartes} \textsuperscript{(82)},\nl
et je montai directement à la \VUgras{salle de billard} \textsuperscript{(74)}.

%%K t13-13-1 p
\VUblancAlinea
13. --- Lorsque la partie fut finie, je redescendis au\nl
rez-de-chaussée et je demandai au garçon une \VUgras{enve}\cc
\VUgras{loppe} \textsuperscript{(84)}, du \VUgras{papier} \textsuperscript{(83)}, un \VUgras{timbre-poste} \textsuperscript{(85)}, une\nl
\VUgras{plume} \textsuperscript{(87)} et de l'\VUgras{encre} \textsuperscript{(86)} pour écrire une lettre.

%%K t13-13-2 p
\VUblancAlinea
Puis j'appelai un \VUgras{cocher} \textsuperscript{(92)} dont la \VUgras{voiture} \textsuperscript{(88)}\nl
était arrêtée en face du café. Je lui demandai le prix\nl
de l'heure et de la course, et, lorsque nous fûmes\nl
d'accord sur le prix, il m'ouvrit la \VUgras{portière} \textsuperscript{(89)} et je\nl
m'installai dans la voiture. Il remonta sur son\nl
\VUgras{siège} \textsuperscript{(91)}, enleva ses chevaux d'un léger coup de\nl
\VUgras{fouet} \textsuperscript{(93)} et nous partîmes faire une délicieuse pro\cc
menade.

\VUsaut{6.0mm}
\VUfilet{22mm}
\end{VUpage}
